% ------------------------------------------------------------
% Page 86
% ------------------------------------------------------------

Lorsque l’on passe sur la route de Meyrueis à Campis on peut voir les ruines de ce qui fut
l’usine des Carburants Forestiers. Les vestiges de 2 fours subsistent encore. Un corps de
bâtiment sert d’écurie pour les chevaux destinés aux promenades. Les jeunes ignorent tout de
cette époque mais nous, les anciens qui l’avons vécue ne pouvons pas oublier.

\vspace{0.8cm}

\begin{flushright}
Le 10 Juillet 1984\\
Robert Cabanel
\end{flushright}

\vfill

\textit{Ce texte a été saisi ( comme on dit maintenant) par Paul Loupiac le 17 Septembre 2002 ;}

\vspace{0.3cm}

\textit{Il m’avait été remis par Georges Rolland en été 2000. Je me suis permis d’ajouter, pour une lecture
plus aisée}

\vspace{0.3cm}

\textit{Les têtes de paragraphes en Gras}

\vspace{0.3cm}

\textit{* Quelques notes --en italique -- qui complètent ou rectifient le témoignage de Robert Cabanel puisque
j’étais à Meyrueis au printemps 1944 et au moulin d’Ayres le 6 Juin ?. Elles montrent combien on
oublie, combien notre mémoire est sélective, et combien certains épisodes nous ont frappé et pas
d’autres. D’où l’importance de témoignages multiples pour les mêmes aventures ou les mêmes
événements.}
